\clearpage\fancyhead[L]{\small COMET.}\hypertarget{native-394-lower}{}\label{native-394-lower}
\begin{center}\textbf{COMET.}\par\vspace{4pt}\hrule\vspace{4pt}\fontsize{10}{12}\selectfont Göttingische gelehrte Anzeigen. No. 149. pp. 1473..1476. 18 September 1815.\par\vspace{4pt}\hrule\vspace{4pt}\end{center}

\par\indent The observations of this most remarkable comet, on which we have reported several times in these pages, were continued by Prof. \textsc{Gauss} at the local observatory for longer than could at first be hoped. Now that the observations have ended, we give them all here together.

\SourceLeaf{395}\begin{center}\fontsize{10}{12}\selectfont\setlength{\tabcolsep}{5pt}\renewcommand{\arraystretch}{1.15}\begin{tabular}{l|c|c|c}
1815 & \shortstack{Mean time\\in Göttingen} & \shortstack{Apparent right\\ascension} & \shortstack{Apparent\\declination}\\\hline
March 20 & $10^{\mathrm u}\;33^{\mathrm m}\;6^{\mathrm s}$ & $54^\circ\;7'\;1''$ & $39^\circ\;7'\;47''\;\mathrm{N.}$\\
21 & $10\;11\;37$ & $54\;33\;21$ & $39\;36\;57$\\
25 & $9\;46\;38$ & $56\;28\;50$ & $41\;38\;5$\\
30 & $9\;50\;57$ & $59\;13\;3$ & $44\;10\;27$\\
April 2 & $9\;11\;55$ & $61\;2\;27$ & $45\;39\;57$\\
June 12 & $10\;53\;31$ & $165\;3\;3$ & $52\;8\;25$\\
30 & $10\;49\;12$ & $184\;39\;24$ & $39\;22\;5$\\
July 13 & $10\;57\;20$ & $194\;32\;52$ & $29\;53\;54$\\
27 & $10\;4\;59$ & $203\;0\;56$ & $20\;37\quad:$\\
29 & $10\;13\;58$ & $204\;6\;3$ & $19\;24\;13$\\
Aug. 4 & $10\;14\;7$ & $207\;14\;8$ & $15\;53\;49$\\
25 & $9\;1\;44$ & $217\;1\;31$ & $5\;33\;36$\\
\end{tabular}\end{center}

\par\indent We should also note that on 13 and 27 July and 4 August, the comparison stars were taken from \textsc{Piazzi}’s catalogue. On all other days, stars were used whose positions were reduced from the \emph{Histoire Céleste} or the \emph{Mémoires de l’Ac. de Paris} 1790. All observations were made with the ring micrometer of the 10-foot \textsc{Herschel} telescope except the one on 12 June, when the distance from no. 219 in the Great Bear in \textsc{Bode}’s numbering and the position angle were measured with an excellent \textsc{Frauenhofer} heliometer.

\par\indent The elliptical orbit that Prof. \textsc{Gauss} determined around the middle of June and that was reported in issue no. 105 of this journal [p. 385 of this volume] has continued to agree well with all the later observations. The later determinations of the orbit by Messrs. \textsc{Bessel} and \textsc{Nicolai}, based on a more careful examination of the observations, also differ little from it. There can now be only a small uncertainty in this comet’s short orbital period. We give these determinations together here.

\Needspace{21\baselineskip}\begin{center}
\emph{Elliptical elements of the comet calculated by Prof.} \textsc{Bessel.}
\end{center}

\begin{center}\fontsize{10}{12}\selectfont Passage through perihelion, 26 April.\par $0^{\mathrm u}\;5^{\mathrm m}\;14^{\mathrm s}5$ Mean time in Paris\end{center}
\begin{center}\fontsize{9}{11}\selectfont\setlength{\tabcolsep}{2pt}\renewcommand{\arraystretch}{1.15}\begin{tabular}{lcr}
Longitude of the node & $.\;.\;.\;.\;.\;.\;.\;.\;.\;.\;.\;.\;.$ & $83^\circ\;28'\;46''18$\\
Longitude of perihelion & $.\;.\;.\;.\;.\;.\;.\;.\;.\;.\;.\;.$ & $149\;2\;29.13$\\
\multicolumn{3}{c}{\shortstack{both for 1 January 1815 and assumed\\fixed with respect to the stars}}\\
Inclination of orbit & $.\;.\;.\;.\;.\;.\;.\;.\;.\;.\;.\;.\;.$ & $44\;29\;53.71$\\
Logarithm of the least distance & $.\;.\;.\;.\;.\;.\;.$ & $0.0837950$\\
Eccentricity & $.\;.\;.\;.\;.\;.\;.\;.\;.\;.\;.\;.\;.\;.\;.$ & $0.93112771$\\
Semimajor axis & $.\;.\;.\;.\;.\;.\;.\;.\;.\;.\;.\;.\;.\;.$ & $17.60964$\\
Sidereal orbital period & $.\;.\;.\;.\;.\;.\;.\;.\;.\;.\;.\;.$ & $73.89682\;\text{years}$\\
\end{tabular}\end{center}

\Needspace{21\baselineskip}\begin{center}
\emph{Elliptical elements of the comet calculated by Mr.} \textsc{Nicolai}, \emph{at the Seeberg Observatory.}
\end{center}

\begin{center}\fontsize{10}{12}\selectfont Passage through perihelion, 26 April.\par $0^{\mathrm u}\;55^{\mathrm m}\;32^{\mathrm s}$ Mean time at the Seeberg Observatory\end{center}
\begin{center}\fontsize{9}{11}\selectfont\setlength{\tabcolsep}{2pt}\renewcommand{\arraystretch}{1.15}\begin{tabular}{lcr}
Longitude of perihelion & $.\;.\;.\;.\;.\;.\;.\;.\;.\;.\;.\;.$ & $149^\circ\;3'\;25''3$\\
Longitude of the node & $.\;.\;.\;.\;.\;.\;.\;.\;.\;.\;.\;.\;.\;.$ & $83\;28\;52.3$\\
\multicolumn{3}{c}{both for 26 April and fixed with respect to the stars}\\
Inclination of orbit & $.\;.\;.\;.\;.\;.\;.\;.\;.\;.\;.\;.\;.$ & $44\;29\;46.0$\\
Logarithm of the least distance & $.\;.\;.\;.\;.\;.\;.$ & $0.0837490$\\
Eccentricity & $.\;.\;.\;.\;.\;.\;.\;.\;.\;.\;.\;.\;.\;.\;.\;.$ & $0.93029341$\\
Semimajor axis & $.\;.\;.\;.\;.\;.\;.\;.\;.\;.\;.\;.\;.\;.$ & $17.39704$\\
Sidereal orbital period & $.\;.\;.\;.\;.\;.\;.\;.\;.\;.\;.\;.$ & $72.564\;\text{Julian years.}$\\
\end{tabular}\end{center}

\Needspace{13\baselineskip}\SourceLeaf{396}\begingroup\setlength{\rightskip}{0pt plus 1em}\par\indent Comparing the latter elements with the most recent observations made here gave the following\par\endgroup

\begin{center}
\fontsize{10}{12}\selectfont\setlength{\tabcolsep}{9pt}\renewcommand{\arraystretch}{1.15}\begin{tabular}{l | c | c}
& \multicolumn{2}{c}{Differences} \\
& Right ascension & Declination \\
\hline
June 30 & $-50''4$ & $+15''9$ \\
July 13 & $-16.6$ & $-44.1$ \\
27 & $-23.8$ & \\
29 & $+3.6$ & $-47.8$ \\
August 4 & $-6.5$ & $-38.1$ \\
25 & $-2.4$ & $-47.5$ \\
\end{tabular}
\end{center}

\Needspace{13\baselineskip}\par\indent We also take this opportunity to report the observations of \emph{Juno} that Prof. \textsc{Gauss} made around the time of opposition. This planet was extremely faint on this occasion, and it does not appear to have been observed at any other observatory.

\begin{center}\fontsize{10}{12}\selectfont\setlength{\tabcolsep}{5pt}\renewcommand{\arraystretch}{1.15}\begin{tabular}{l|c|c|c}
1815 & \shortstack{Mean time\\in Göttingen} & \shortstack{Apparent right\\ascension} & \shortstack{Apparent\\declination}\\\hline
March 1 & $10^{\mathrm u}\;22^{\mathrm m}\;0''5$ & $197^\circ\;7'\;56''7$ & $2^\circ\;19'\;3''0\;\mathrm{S.}$\\
29 & $10\;8\;57.0$ & $192\;31\;51.8$ & $1\;32\;30.1\;\mathrm{N.}$\\
April 8 & $9\;27\;42.9$ & $190\;34\;51.9$ & $2\;53\;24.2\;\mathrm{N.}$\\
\end{tabular}\end{center}

\Needspace{14\baselineskip}\par\indent Mr. \textsc{Nicolai} undertook to compare these observations with the latest elements calculated by Mr. \textsc{Möbius} (see Mon. Corresp. 1813, vol. XXVIII, p. 577). After increasing the epoch by $4'$ $55''$, he found the following

\begin{center}
\fontsize{10}{12}\selectfont\setlength{\tabcolsep}{6pt}\renewcommand{\arraystretch}{1.15}\begin{tabular}{l | c | c | c | c}
& \multicolumn{4}{c}{Differences} \\
& Right ascension & Declination & Longitude & Latitude \\
\hline
March 1 & $-20''6$ & $-49.7$ & $0$ & $-53''7$ \\
29 & $-12.2$ & $-50.5$ & $+8.6$ & $-51.2$ \\
April 8 & $-16.3$ & $-45.1$ & $+2.9$ & $-47.8$ \\
\end{tabular}
\end{center}

\noindent and from these the opposition with the Sun

\begin{center}31 March 1815.\quad $12^{\mathrm u}\;56^{\mathrm m}\;43^{\mathrm s}$ Mean time in Göttingen\end{center}
\begin{center}\fontsize{9}{11}\selectfont\setlength{\tabcolsep}{2pt}\renewcommand{\arraystretch}{1.15}\begin{tabular}{lcr}
True longitude & $.\;.\;.\;.\;.\;.\;.\;.\;.\;.\;.\;.\;.$ & $190^\circ\;25'\;6''6$\\
Northern geocentric latitude & $.\;.\;.\;.\;.\;.\;.$ & $6\;29\;10.0$\\
\end{tabular}\end{center}

\Needspace{21\baselineskip}\par\indent By combining this opposition with the three preceding ones, this skilled astronomer also obtained the following

\Needspace{18\baselineskip}\begin{center}
\emph{New elements of Juno.}
\end{center}

\begin{center}\fontsize{9}{11}\selectfont\setlength{\tabcolsep}{2pt}\renewcommand{\arraystretch}{1.15}\begin{tabular}{lcr}
Epoch of the mean longitude 1816 & $.\;.$ & $230^\circ\;11'\;34''2$\\
\multicolumn{3}{c}{for the meridian of Göttingen}\\
Daily tropical motion & $.\;.\;.$ & $812''9304$\\
Longitude of perihelion 1816 & $.\;.\;.\;.\;.$ & $53^\circ\;14'\;53''8$\\
Longitude of the ascending node 1816 &  & $171\;9\;58.9$\\
Inclination of orbit & $.\;.\;.\;.\;.\;.\;.$ & $13\;14\;0.1$\\
Eccentricity angle & $.\;.\;.\;.\;.\;.$ & $14\;43\;28.84$\\
Logarithm of the semimajor axis & $.$ & $0.4266844$\\
\end{tabular}\end{center}

\par\indent An ephemeris calculated by Mr. \textsc{Nicolai} from these elements will be published in the astronomical yearbook for 1818.
